\section{Introduction}\label{sec:introduction}

As established in \cite{happelSheavesMinecraft2021}, information that depends
on location can be modelled using \emph{sheaves}. We wish to expand on this idea
in the context of the popular video game \emph{Minecraft}. Its world is a subset
of $\RR^3$, inhabited by two kinds of objects:
\begin{enumerate}[label=(\arabic*)]
	\item \emph{blocks}, which are unit cubes placed on the lattice
	      $\ZZ^3 \subset \RR^3$, and
	\item \emph{entities}, such as dropped items, players and enemies. Several
	      entities may occupy the same location at the same time.
\end{enumerate}

To render a (more or less) infinite world in finite time, Minecraft
approximates it by finitely many pieces, a technique known as
\emph{chunking}. A \emph{chunk} is a bounded region of $\RR^3$, together with
the information about which blocks and which entities are placed inside of it.
In \cite{happelSheavesMinecraft2021} we saw that chunking can in fact be
modelled using (pre)sheaves. To this end, fix a width $w$, a depth $d$ and a
height $h$, all positive integers, and equip $\RR^3$ with the topology $\tau$
generated by the closed boxes
\[
	C_{i,j,k}^{w,d,h} = [i, i+w] \times [j, j+d] \times [k, k+h],
	\qquad (i,j,k) \in w\ZZ \times d\ZZ \times h\ZZ.
\]
These boxes tile the space $\RR^3$, as seen in \cref{fig:chunk-tiling}. Since
the boxes are closed, two neighbouring chunks share a common face, so $\tau$
contains not only (unions of) chunks but also their faces, edges and vertices
as open sets.

\begin{figure}[ht]
	\centering
	\begin{tikzpicture}[
			x={(1.0cm,0cm)}, y={(0.55cm,0.38cm)}, z={(0cm,1.0cm)},
			scale=1.15, line join=round, line cap=round,
			grid/.style={line width=0.5pt},
			stub/.style={line width=0.4pt, dash pattern=on 1.5pt off 1.5pt},
		]
		% chunk dimensions (in picture units) and number of chunks shown
		\def\w{1}\def\d{1}\def\h{1.25}
		\def\nx{4}\def\ny{3}\def\nz{2}
		\pgfmathsetmacro{\X}{\nx*\w}
		\pgfmathsetmacro{\Y}{\ny*\d}
		\pgfmathsetmacro{\Z}{\nz*\h}

		% faces of the block, shaded by ordered dithering (black/white only)
		\fill[white] (0,0,0) -- (\X,0,0) -- (\X,0,\Z) -- (0,0,\Z) -- cycle;
		\fill[pattern=dots] (0,0,\Z) -- (\X,0,\Z) -- (\X,\Y,\Z) -- (0,\Y,\Z) -- cycle;
		\fill[pattern=crosshatch dots] (\X,0,0) -- (\X,\Y,0) -- (\X,\Y,\Z) -- (\X,0,\Z) -- cycle;

		% one distinguished chunk C_{0,0,h}^{w,d,h} on the front face
		\fill[pattern=north east lines] (\w,0,\h) -- (2*\w,0,\h) -- (2*\w,0,2*\h) -- (\w,0,2*\h) -- cycle;

		% grid lines on the three visible faces
		\foreach \i in {0,...,\nx} {
				\draw[grid] (\i*\w,0,0) -- (\i*\w,0,\Z) -- (\i*\w,\Y,\Z);
				\draw[stub] (\i*\w,0,0) -- (\i*\w,0,-0.45);
				\draw[stub] (\i*\w,\Y,\Z) -- (\i*\w,\Y+0.8,\Z);
			}
		\foreach \k in {0,...,\nz} {
				\draw[grid] (0,0,\k*\h) -- (\X,0,\k*\h) -- (\X,\Y,\k*\h);
				\draw[stub] (0,0,\k*\h) -- (-0.45,0,\k*\h);
				\draw[stub] (\X,\Y,\k*\h) -- (\X,\Y+0.8,\k*\h);
			}
		\foreach \j in {0,...,\ny} {
				\draw[grid] (0,\j*\d,\Z) -- (\X,\j*\d,\Z) -- (\X,\j*\d,0);
				\draw[stub] (0,\j*\d,\Z) -- (-0.45,\j*\d,\Z);
				\draw[stub] (\X,\j*\d,0) -- (\X,\j*\d,-0.45);
			}

		% outline
		\draw[line width=0.9pt] (0,0,0) -- (\X,0,0) -- (\X,\Y,0) -- (\X,\Y,\Z)
		-- (0,\Y,\Z) -- (0,0,\Z) -- cycle;
		\draw[line width=0.9pt] (0,0,\Z) -- (\X,0,\Z) -- (\X,\Y,\Z);
		\draw[line width=0.9pt] (\X,0,0) -- (\X,0,\Z);

		% dimensions of a single chunk
		\draw[{Bar[width=4pt]}-{Bar[width=4pt]}, line width=0.4pt]
		(0,0,-0.75) -- node[fill=white, inner sep=1.5pt] {$w$} (\w,0,-0.75);
		\draw[{Bar[width=4pt]}-{Bar[width=4pt]}, line width=0.4pt]
		(\X+0.55,0,-0.1) -- node[below right=-1pt, inner sep=1.5pt] {$d$} (\X+0.55,\d,-0.1);
		\draw[{Bar[width=4pt]}-{Bar[width=4pt]}, line width=0.4pt]
		(-0.85,0,0) -- node[fill=white, inner sep=1.5pt] {$h$} (-0.85,0,\h);

		% label of the distinguished chunk
		\node[fill=white, inner sep=1.5pt, font=\small] at (1.5*\w,0,1.5*\h) {$C_{w,0,h}^{w,d,h}$};
	\end{tikzpicture}
	\caption{The boxes $C_{i,j,k}^{w,d,h}$ tile the space $\RR^3$. Only finitely many
		chunks are drawn. The dashed lines indicate that the tiling continues
		in every direction.}
	\label{fig:chunk-tiling}
\end{figure}

A sheaf over chunks is then given by a contravariant functor
$P \colon \tau^{\mathrm{op}} \to \Set$, which assigns to every chunk the set
\[
	P(C_{i,j,k}^{w,d,h}) = \Blocks_{w,d,h} \times \Entities_{w,d,h}
\]
of its possible contents, where
\begin{align*}
	\Blocks_{w,d,h}   & \coloneqq \Hom_{\Set}(\NN_{<w} \times \NN_{<d} \times \NN_{<h},\ \Block), \\
	\Entities_{w,d,h} & \coloneqq \Hom_{\Set}([0,w) \times [0,d) \times [0,h),\ \List(\Entity)),
\end{align*}
where $\Block$ is the set of all possible blocks (say
$\mathrm{AirBlock} \in \Block$, $\mathrm{DirtBlock} \in \Block$,
$\mathrm{WaterBlock} \in \Block$, and so on), and $\Entity$ is the set of all
possible entities (say $\mathrm{Player} \in \Entity$,
$\mathrm{Zombie} \in \Entity$, and so on). A function in $\Blocks_{w,d,h}$
assigns a block to every position of a chunk, and a function in
$\Entities_{w,d,h}$ assigns to every point of a chunk the entities located
there. Here $\List(A)$ is the inductive set with
$\List(A) \cong A \times \List(A) + 1$. Since the value at a point is a list,
multiple entities may share a position.

Crucially, $P$ needs to satisfy the \emph{sheaf conditions}, so for every open
set $U$ and every open cover $U = \bigcup_{\iota \in I} U_\iota$,
\begin{enumerate}[label=(\roman*)]
	\item \emph{locality}: if $s, t \in P(U)$ satisfy
	      $\restr{s}{U_\iota} = \restr{t}{U_\iota}$ for all $\iota \in I$, then
	      $s = t$, and
	\item \emph{gluing}: if $s_\iota \in P(U_\iota)$ is a family with
	      $\restr{s_\iota}{U_\iota \cap U_\kappa} = \restr{s_\kappa}{U_\iota \cap U_\kappa}$
	      for all $\iota, \kappa \in I$, then there is some $s \in P(U)$ with
	      $\restr{s}{U_\iota} = s_\iota$ for all $\iota \in I$.
\end{enumerate}
These conditions are important because for the boundaries between chunks they
ensure that two neighbouring chunks, which share a common face, fit together.
For further information about the need for these conditions, consider
\cite{happelSheavesMinecraft2021}.

We may also obtain $P$ as the product $P = E \times B$ of an \emph{entity
	sheaf} $E$ and a \emph{block sheaf} $B$, with $E(C_{i,j,k}^{w,d,h}) = \Entities_{w,d,h}$
and $B(C_{i,j,k}^{w,d,h}) = \Blocks_{w,d,h}$. This works because a functor into a product
of categories is the same as a pair of functors,
\[
	\functorcat{\tau^{\mathrm{op}}}{\Set \times \Set}
	\cong \functorcat{\tau^{\mathrm{op}}}{\Set} \times \functorcat{\tau^{\mathrm{op}}}{\Set},
\]
\todo[inline]{Is this really true for sheaves?}
and $P$ is the image of the pair $(B, E)$ under this isomorphism, composed with
the product functor ${\times} \colon \Set \times \Set \to \Set$.

\begin{figure}[ht]
	\centering
	% ----------------------------------------------------------------
	% Shared drawing code (local to this figure).
	% Same parallel projection as in Fig.~\ref{fig:chunk-tiling}.
	% ----------------------------------------------------------------
	\tikzset{
		sheafpic/.style={
				x={(1.0cm,0cm)}, y={(0.55cm,0.38cm)}, z={(0cm,1.0cm)},
				scale=0.98, line join=round, line cap=round},
		chunkline/.style={line width=0.4pt, dash pattern=on 2pt off 1.5pt},
		chunkstub/.style={line width=0.3pt, dash pattern=on 0.8pt off 1.6pt},
		vox/.style={line width=0.3pt, preaction={fill=white}},
		% ordered "dithering": black and white only
		dither sparse/.style={pattern={Dots[distance=2.2pt, radius=0.28pt]}},
		dither mid/.style={pattern={Dots[distance=1.3pt, radius=0.3pt]}},
		dither lines/.style={pattern={Lines[angle=45, distance=1.3pt, line width=0.3pt]}},
		dither vert/.style={pattern={Lines[angle=90, distance=1.1pt, line width=0.3pt]}},
		dither dark/.style={pattern={Hatch[angle=45, distance=1.1pt, line width=0.3pt]}},
	}
	% a unit cube at voxel (#1,#2,#3), faces styled #4 (top), #5 (front), #6 (right)
	\newcommand{\vox}[6]{%
		\draw[vox,#5] (#1,#2,#3) -- ++(1,0,0) -- ++(0,0,1) -- ++(-1,0,0) -- cycle;
		\draw[vox,#6] ({#1+1},#2,#3) -- ++(0,1,0) -- ++(0,0,1) -- ++(0,-1,0) -- cycle;
		\draw[vox,#4] (#1,#2,{#3+1}) -- ++(1,0,0) -- ++(0,1,0) -- ++(-1,0,0) -- cycle;}
	\newcommand{\dirt}[3]{\vox{#1}{#2}{#3}{dither sparse}{dither mid}{dither dark}}
	\newcommand{\grass}[3]{\vox{#1}{#2}{#3}{dither sparse}{dither mid}{dither dark}}
	\newcommand{\wood}[3]{\vox{#1}{#2}{#3}{dither mid}{dither vert}{dither dark}}
	\newcommand{\leaves}[3]{\vox{#1}{#2}{#3}{dither sparse}{dither lines}{dither dark}}
	\newcommand{\chest}[3]{%
		\begin{scope}[shift={({#1+0.1},{#2+0.1},#3)}, scale=0.8]
			\vox{0}{0}{0}{}{}{dither sparse}
			\draw[line width=0.3pt] (0,0,0.62) -- (1,0,0.62) -- (1,1,0.62);
			\fill (0.42,0,0.48) rectangle (0.58,0,0.72);
		\end{scope}}
	% the chunk grid of Fig.~\ref{fig:chunk-tiling}, outlined by dashed lines only
	\newcommand{\chunkgrid}{%
		\def\w{1}\def\d{1}\def\h{1.25}
		\def\nx{4}\def\ny{3}\def\nz{2}
		\pgfmathsetmacro{\X}{\nx*\w}
		\pgfmathsetmacro{\Y}{\ny*\d}
		\pgfmathsetmacro{\Z}{\nz*\h}
		\foreach \i in {0,...,\nx} {
				\draw[chunkline] (\i*\w,0,0) -- (\i*\w,0,\Z) -- (\i*\w,\Y,\Z);
				\draw[chunkstub] (\i*\w,0,0) -- (\i*\w,0,-0.45);
				\draw[chunkstub] (\i*\w,\Y,\Z) -- (\i*\w,\Y+0.6,\Z);
			}
		\foreach \k in {0,...,\nz} {
				\draw[chunkline] (0,0,\k*\h) -- (\X,0,\k*\h) -- (\X,\Y,\k*\h);
				\draw[chunkstub] (0,0,\k*\h) -- (-0.45,0,\k*\h);
				\draw[chunkstub] (\X,\Y,\k*\h) -- (\X,\Y+0.6,\k*\h);
			}
		\foreach \j in {0,...,\ny} {
				\draw[chunkline] (0,\j*\d,\Z) -- (\X,\j*\d,\Z) -- (\X,\j*\d,0);
				\draw[chunkstub] (0,\j*\d,\Z) -- (-0.45,\j*\d,\Z);
				\draw[chunkstub] (\X,\j*\d,0) -- (\X,\j*\d,-0.45);
			}
		% hidden edges of the distinguished chunk C_{w,0,h}^{w,d,h}
		\draw[chunkline] (\w,\d,\h) -- (2*\w,\d,\h);
		\draw[chunkline] (\w,0,\h) -- (\w,\d,\h) -- (\w,\d,2*\h);
		\draw[chunkline] (2*\w,0,\h) -- (2*\w,\d,\h) -- (2*\w,\d,2*\h);
		% dimensions of a single chunk
		\draw[{Bar[width=4pt]}-{Bar[width=4pt]}, line width=0.4pt]
		(0,0,-0.75) -- node[fill=white, inner sep=1.5pt] {$w$} (\w,0,-0.75);
		\draw[{Bar[width=4pt]}-{Bar[width=4pt]}, line width=0.4pt]
		(\X+0.55,0,-0.1) -- node[below right=-1pt, inner sep=1.5pt] {$d$} (\X+0.55,\d,-0.1);
		\draw[{Bar[width=4pt]}-{Bar[width=4pt]}, line width=0.4pt]
		(-0.85,0,0) -- node[fill=white, inner sep=1.5pt] {$h$} (-0.85,0,\h);
	}
	% voxel coordinates inside C_{w,0,h}^{w,d,h}: 5 x 5 x 6 cubes of side 0.2
	\tikzset{inchunk/.style={shift={(1.1,0,1.25+0.025)}, scale=0.2}}
	\begin{subfigure}[b]{0.49\textwidth}
		\centering
		\begin{tikzpicture}[sheafpic]
			\chunkgrid
			\begin{scope}[inchunk]
				\leaves{0}{4}{4}
\leaves{1}{4}{4}
\leaves{1}{4}{5}
\leaves{2}{4}{4}
\dirt{0}{3}{0}
\grass{0}{3}{1}
\leaves{0}{3}{4}
\leaves{0}{3}{5}
\dirt{1}{3}{0}
\grass{1}{3}{1}
\wood{1}{3}{2}
\wood{1}{3}{3}
\leaves{1}{3}{4}
\leaves{1}{3}{5}
\leaves{2}{3}{4}
\leaves{2}{3}{5}
\dirt{0}{2}{0}
\grass{0}{2}{1}
\leaves{0}{2}{4}
\dirt{1}{2}{0}
\grass{1}{2}{1}
\leaves{1}{2}{4}
\leaves{1}{2}{5}
\leaves{2}{2}{4}
\dirt{0}{1}{0}
\grass{0}{1}{1}
\dirt{1}{1}{0}
\grass{1}{1}{1}
\dirt{2}{1}{0}
\grass{2}{1}{1}
\dirt{3}{1}{0}
\grass{3}{1}{1}
\dirt{0}{0}{0}
\grass{0}{0}{1}
\dirt{1}{0}{0}
\grass{1}{0}{1}
\dirt{2}{0}{0}
\grass{2}{0}{1}
\dirt{3}{0}{0}
\grass{3}{0}{1}
\chest{3}{0}{2}
			\end{scope}
			% front edges of C_{w,0,h}^{w,d,h} lie in front of the island
			\draw[chunkline] (1,0,1.25) rectangle (2,0,2.5);
		\end{tikzpicture}
		\caption{The block sheaf $B$: blocks placed in the chunk $C_{w,0,h}^{w,d,h}$.}
		\label{fig:block-sheaf}
	\end{subfigure}%
	\hfill
	\begin{subfigure}[b]{0.49\textwidth}
		\centering
		\begin{tikzpicture}[sheafpic]
			\chunkgrid
			\begin{scope}[inchunk]
				% a player standing where the island's grass would be
				\fill (2.6,0.5,2) circle[radius=0.12];
				\begin{scope}[shift={(2.6,0.5,2)}, line width=0.5pt]
					\draw (-0.3cm,0cm) -- (0cm,0.8cm) -- (0.3cm,0cm);
					\draw (0cm,0.8cm) -- (0cm,1.45cm);
					\draw (-0.45cm,0.9cm) -- (0cm,1.3cm) -- (0.45cm,0.9cm);
					\draw[preaction={fill=white}] (0cm,1.67cm) circle[radius=0.22cm];
				\end{scope}
			\end{scope}
		\end{tikzpicture}
		\caption{The entity sheaf $E$: a player standing in the chunk $C_{w,0,h}^{w,d,h}$.}
		\label{fig:entity-sheaf}
	\end{subfigure}
	\caption{The two components of $P = E \times B$ over the same chunk
		$C_{w,0,h}^{w,d,h}$. Blocks (left) sit on the lattice $\ZZ^3$, entities (right)
		at arbitrary points.}
	\label{fig:block-entity-sheaf}
\end{figure}

The topology $\tau$ induces a labelled graph $\gamma$. For each chunk, we add a
node labelled by \cubeglyph, and for their faces we add nodes labelled by
\faceglyph, and similarly for their edges \edgeglyph{} and vertices
\vertexglyph. These nodes are then connected by edges iff one of the two
corresponding cells lies on the boundary of the other, where each edge needs to
join cells of consecutive dimensions, that is, a chunk \cubeglyph{} with a face
\faceglyph, a face \faceglyph{} with an edge \edgeglyph, or an edge
\edgeglyph{} with a vertex \vertexglyph.

We call a sheaf $P \colon \tau^{\mathrm{op}} \to \Set$ with
$P(C_{i,j,k}^{w,d,h}) = \Blocks_{w,d,h} \times \Entities_{w,d,h}$ a \emph{Euclidean
	Mathcraft sheaf}. This would yield the induced labelled graph $X$ in
\cref{fig:chunk-lattice}. We could also consider other topologies which also
yield such a labelled graph, like one that induces the graph $Y$ as shown in
\cref{fig:wormhole-lattice}, where we added a \enquote{wormhole}, which is
given by an additional face \faceglyph{} shared by two chunks \cubeglyph. This
justifies the term \emph{Euclidean} to distinguish between the standard
topology and others.

\begin{figure}[ht]
	\centering
	\tikzset{
		latticepic/.style={
				x={(1.0cm,0cm)}, y={(0.55cm,0.38cm)}, z={(0cm,1.0cm)},
				scale=1.3, line join=round, line cap=round},
		celledge/.style={line width=0.5pt},
		hasseedge/.style={line width=0.3pt},
		latticestub/.style={line width=0.3pt, dash pattern=on 0.8pt off 1.6pt},
		cellglyph/.style={line width=0.45pt, preaction={fill=white}},
		wormhole/.style={line width=0.7pt},
	}
	% glyph size (in units of the chunk width)
	\def\gs{0.2}
	% node glyphs: every cell is drawn as a small copy of itself
	\newcommand{\vertexnode}[3]{\draw[cellglyph] (#1,#2,#3) circle[radius=1.6pt];}
	% edge at (#1,#2,#3) along the axis (#4,#5,#6)
	\newcommand{\edgenode}[6]{%
		\draw[line width=2.6pt, line cap=round]
		({#1-0.5*\gs*#4},{#2-0.5*\gs*#5},{#3-0.5*\gs*#6}) -- ++({\gs*#4},{\gs*#5},{\gs*#6});
		\draw[line width=1.4pt, line cap=round, white]
		({#1-0.5*\gs*#4},{#2-0.5*\gs*#5},{#3-0.5*\gs*#6}) -- ++({\gs*#4},{\gs*#5},{\gs*#6});}
	% face at (#1,#2,#3) spanned by the axes (#4,#5,#6) and (#7,#8,#9)
	\newcommand{\facenode}[9]{%
		\draw[cellglyph]
		({#1-0.5*\gs*(#4+#7)},{#2-0.5*\gs*(#5+#8)},{#3-0.5*\gs*(#6+#9)})
		-- ++({\gs*#4},{\gs*#5},{\gs*#6}) -- ++({\gs*#7},{\gs*#8},{\gs*#9})
		-- ++({-\gs*#4},{-\gs*#5},{-\gs*#6}) -- cycle;}
	\newcommand{\chunknode}[3]{%
		\begin{scope}[shift={({#1-0.5*\gs},{#2-0.5*\gs},{#3-0.5*\gs})}, scale=\gs]
			\draw[cellglyph] (0,0,0) -- (1,0,0) -- (1,0,1) -- (0,0,1) -- cycle;
			\draw[cellglyph, pattern={Lines[angle=90, distance=1.1pt, line width=0.3pt]}]
			(1,0,0) -- (1,1,0) -- (1,1,1) -- (1,0,1) -- cycle;
			\draw[cellglyph, pattern={Dots[distance=1.2pt, radius=0.3pt]}]
			(0,0,1) -- (1,0,1) -- (1,1,1) -- (0,1,1) -- cycle;
		\end{scope}}
	% the cell complex of a 2 x 2 x 1 block of chunks of height \h, raised by #1.
	% Points of the half-step grid are the centres of the cells, and lines of
	% the half-step grid join every cell to the cells on its boundary.
	\newcommand{\cellcomplex}[1]{%
		\def\h{1.25}
		% lines along x, y and z
		\foreach \B in {0,...,4} {\foreach \C in {0,...,2} {
						\pgfmathsetmacro{\yy}{\B/2}\pgfmathsetmacro{\zz}{#1+\C/2*\h}
						\ifodd\B \draw[hasseedge] (0,\yy,\zz) -- (2,\yy,\zz);
						\else\ifodd\C \draw[hasseedge] (0,\yy,\zz) -- (2,\yy,\zz);
							\else
								\draw[latticestub] (-0.35,\yy,\zz) -- (2.35,\yy,\zz);
								\draw[celledge] (0,\yy,\zz) -- (2,\yy,\zz);
							\fi\fi}}
		\foreach \A in {0,...,4} {\foreach \C in {0,...,2} {
						\pgfmathsetmacro{\xx}{\A/2}\pgfmathsetmacro{\zz}{#1+\C/2*\h}
						\ifodd\A \draw[hasseedge] (\xx,0,\zz) -- (\xx,2,\zz);
						\else\ifodd\C \draw[hasseedge] (\xx,0,\zz) -- (\xx,2,\zz);
							\else
								\draw[latticestub] (\xx,-0.35,\zz) -- (\xx,2.35,\zz);
								\draw[celledge] (\xx,0,\zz) -- (\xx,2,\zz);
							\fi\fi}}
		\foreach \A in {0,...,4} {\foreach \B in {0,...,4} {
						\pgfmathsetmacro{\xx}{\A/2}\pgfmathsetmacro{\yy}{\B/2}
						\ifodd\A \draw[hasseedge] (\xx,\yy,#1) -- (\xx,\yy,#1+\h);
						\else\ifodd\B \draw[hasseedge] (\xx,\yy,#1) -- (\xx,\yy,#1+\h);
							\else
								\draw[latticestub] (\xx,\yy,#1-0.3) -- (\xx,\yy,#1+\h+0.3);
								\draw[celledge] (\xx,\yy,#1) -- (\xx,\yy,#1+\h);
							\fi\fi}}
		% nodes, from back to front
		\foreach \B in {4,...,0} {\foreach \A in {0,...,4} {\foreach \C in {0,...,2} {
							\pgfmathsetmacro{\xx}{\A/2}\pgfmathsetmacro{\yy}{\B/2}
							\pgfmathsetmacro{\zz}{#1+\C/2*\h}
							\ifodd\A
								\ifodd\B
									\ifodd\C \chunknode{\xx}{\yy}{\zz}\else \facenode{\xx}{\yy}{\zz}{1}{0}{0}{0}{1}{0}\fi
								\else
									\ifodd\C \facenode{\xx}{\yy}{\zz}{1}{0}{0}{0}{0}{1}\else \edgenode{\xx}{\yy}{\zz}{1}{0}{0}\fi
								\fi
							\else
								\ifodd\B
									\ifodd\C \facenode{\xx}{\yy}{\zz}{0}{1}{0}{0}{0}{1}\else \edgenode{\xx}{\yy}{\zz}{0}{1}{0}\fi
								\else
									\ifodd\C \edgenode{\xx}{\yy}{\zz}{0}{0}{1}\else \vertexnode{\xx}{\yy}{\zz}\fi
								\fi
							\fi}}}
	}
	\begin{subfigure}[b]{0.49\textwidth}
		\centering
		\begin{tikzpicture}[latticepic]
			\cellcomplex{0}
		\end{tikzpicture}
		\caption{The graph $X$.}
		\label{fig:chunk-lattice}
	\end{subfigure}%
	\hfill
	\begin{subfigure}[b]{0.49\textwidth}
		\centering
		\begin{tikzpicture}[latticepic]
			\cellcomplex{2.6}
			\cellcomplex{0}
			% a wormhole: an additional face shared by two chunks
			\draw[wormhole] (0.5,0.5,3.225) to[out=180, in=90] (-0.9,0.5,1.925);
			\draw[wormhole] (-0.9,0.5,1.925) to[out=-90, in=180] (0.5,0.5,0.625);
			\facenode{-0.9}{0.5}{1.925}{0}{1}{0}{0}{0}{1}
			\chunknode{0.5}{0.5}{3.225}
			\chunknode{0.5}{0.5}{0.625}
		\end{tikzpicture}
		\caption{The graph $Y$, with a wormhole face.}
		\label{fig:wormhole-lattice}
	\end{subfigure}
	\caption{Graphs underlying Mathcraft sheaves, drawn for $2 \times 2 \times 1$
		chunks. Each node stands for a cell of $\tau$: vertices are drawn as circles, edges as bars, faces as squares
		and chunks as cubes. Lines join every cell to the cells on its boundary.}
	\label{fig:lattices}
\end{figure}

In fact, in Minecraft this is exactly what is used for portals, such as
Nether portals or End portals. This becomes apparent when considering a block
sheaf over $Y$, as shown in \cref{fig:nether-portal}.

\begin{figure}[ht]
	\centering
	\tikzset{
		portalpic/.style={
				x={(1.0cm,0cm)}, y={(0.55cm,0.38cm)}, z={(0cm,1.0cm)},
				scale=1.55, line join=round, line cap=round},
		chunkline/.style={line width=0.4pt, dash pattern=on 2pt off 1.5pt},
		chunkstub/.style={line width=0.3pt, dash pattern=on 0.8pt off 1.6pt},
		vox/.style={line width=0.3pt, preaction={fill=white}},
		% ordered "dithering": black and white only
		dither sparse/.style={pattern={Dots[distance=2.2pt, radius=0.28pt]}},
		dither mid/.style={pattern={Dots[distance=1.3pt, radius=0.3pt]}},
		dither lines/.style={pattern={Lines[angle=45, distance=1.3pt, line width=0.3pt]}},
		dither vert/.style={pattern={Lines[angle=90, distance=1.1pt, line width=0.3pt]}},
		dither horiz/.style={pattern={Lines[angle=0, distance=1.1pt, line width=0.3pt]}},
		dither rack/.style={pattern={Lines[angle=-45, distance=1.2pt, line width=0.3pt]}},
		dither dark/.style={pattern={Hatch[angle=45, distance=1.1pt, line width=0.3pt]}},
		dither swirl/.style={pattern={Lines[angle=60, distance=1.5pt, line width=0.3pt]}},
		% obsidian: black with white speckles and white seams
		obsidian/.style={line width=0.3pt, draw=white, fill=black,
				postaction={pattern={Dots[distance=#1, radius=0.3pt]}, pattern color=white}},
		cellglyph/.style={line width=0.45pt, preaction={fill=white}},
		wormhole/.style={line width=0.7pt},
	}
	% a unit cube at voxel (#1,#2,#3), faces styled #4 (top), #5 (front), #6 (right)
	\newcommand{\vox}[6]{%
		\draw[vox,#5] (#1,#2,#3) -- ++(1,0,0) -- ++(0,0,1) -- ++(-1,0,0) -- cycle;
		\draw[vox,#6] ({#1+1},#2,#3) -- ++(0,1,0) -- ++(0,0,1) -- ++(0,-1,0) -- cycle;
		\draw[vox,#4] (#1,#2,{#3+1}) -- ++(1,0,0) -- ++(0,1,0) -- ++(-1,0,0) -- cycle;}
	\newcommand{\dirt}[3]{\vox{#1}{#2}{#3}{dither sparse}{dither mid}{dither dark}}
	\newcommand{\grass}[3]{\vox{#1}{#2}{#3}{dither sparse}{dither mid}{dither dark}}
	\newcommand{\wood}[3]{\vox{#1}{#2}{#3}{dither mid}{dither vert}{dither dark}}
	\newcommand{\leaves}[3]{\vox{#1}{#2}{#3}{dither sparse}{dither lines}{dither dark}}
	\newcommand{\rack}[3]{\vox{#1}{#2}{#3}{dither mid}{dither rack}{dither dark}}
	\newcommand{\lava}[3]{\vox{#1}{#2}{#3}{dither horiz}{dither horiz}{dither dark}}
	\newcommand{\obs}[3]{%
		\draw[obsidian=1.6pt] (#1,#2,#3) -- ++(1,0,0) -- ++(0,0,1) -- ++(-1,0,0) -- cycle;
		\draw[obsidian=2.4pt] ({#1+1},#2,#3) -- ++(0,1,0) -- ++(0,0,1) -- ++(0,-1,0) -- cycle;
		\draw[obsidian=1.1pt] (#1,#2,{#3+1}) -- ++(1,0,0) -- ++(0,1,0) -- ++(-1,0,0) -- cycle;}
	% the portal itself is a flat pane in the front plane of the frame
	\newcommand{\portal}[3]{%
		\fill[white] (#1,#2,#3) -- ++(1,0,0) -- ++(0,0,1) -- ++(-1,0,0) -- cycle;
		\fill[dither swirl] (#1,#2,#3) -- ++(1,0,0) -- ++(0,0,1) -- ++(-1,0,0) -- cycle;}
	\newcommand{\chest}[3]{%
		\begin{scope}[shift={({#1+0.1},{#2+0.1},#3)}, scale=0.8]
			\vox{0}{0}{0}{}{}{dither sparse}
			\draw[line width=0.3pt] (0,0,0.62) -- (1,0,0.62) -- (1,1,0.62);
			\fill (0.42,0,0.48) rectangle (0.58,0,0.72);
		\end{scope}}
	% dashed outlines of a 2 x 2 x 1 block of chunks, raised by #1
	\newcommand{\chunkblock}[1]{%
		\def\h{1.25}
		\foreach \i in {0,1,2} {
				\draw[chunkline] (\i,0,#1) -- (\i,0,#1+\h) -- (\i,2,#1+\h);
				\draw[chunkstub] (\i,0,#1) -- (\i,0,#1-0.3);
				\draw[chunkstub] (\i,2,#1+\h) -- (\i,2.4,#1+\h);
			}
		\foreach \k in {0,1} {
				\draw[chunkline] (0,0,#1+\k*\h) -- (2,0,#1+\k*\h);
				\draw[chunkstub] (0,0,#1+\k*\h) -- (-0.3,0,#1+\k*\h);
				\draw[chunkstub] (2,2,#1+\k*\h) -- (2,2.4,#1+\k*\h);
			}
		\foreach \j in {0,1,2} {
				\draw[chunkline] (0,\j,#1+\h) -- (2,\j,#1+\h) -- (2,\j,#1);
				\draw[chunkstub] (0,\j,#1+\h) -- (-0.3,\j,#1+\h);
				\draw[chunkstub] (2,\j,#1) -- (2.4,\j,#1);
			}
		\draw[chunkline] (2,0,#1) -- (2,2,#1);
		% hidden edges
		\draw[chunkline] (0,0,#1) -- (0,2,#1) -- (2,2,#1);
		\draw[chunkline] (0,2,#1) -- (0,2,#1+\h);
		\draw[chunkline] (1,0,#1) -- (1,2,#1) -- (1,2,#1+\h);
		\draw[chunkline] (0,1,#1) -- (2,1,#1);
		\draw[chunkline] (0,1,#1) -- (0,1,#1+\h);
		\draw[chunkline] (1,1,#1) -- (1,1,#1+\h);
	}
	% voxels of side 0.2, shifted half a voxel to the right within each chunk
	\tikzset{voxels/.style={shift={(0.1,0,#1+0.025)}, scale=0.2}}
	\begin{tikzpicture}[portalpic]
		% the Nether
		\chunkblock{0}
		\begin{scope}[voxels=0]
			\rack{5}{8}{0}\rack{6}{8}{0}\rack{7}{8}{0}\rack{8}{8}{0}\rack{1}{7}{2}
				\rack{1}{7}{3}\rack{2}{7}{0}\rack{2}{7}{1}\rack{5}{7}{0}\lava{6}{7}{0}
				\lava{7}{7}{0}\rack{8}{7}{0}\rack{1}{6}{0}\rack{1}{6}{1}\rack{2}{6}{0}
				\rack{2}{6}{1}\rack{5}{6}{0}\lava{6}{6}{0}\lava{7}{6}{0}\rack{8}{6}{0}
				\rack{5}{5}{0}\rack{6}{5}{0}\rack{7}{5}{0}\rack{8}{5}{0}\rack{5}{3}{0}
				\rack{6}{3}{1}\rack{7}{3}{2}\rack{8}{3}{0}\rack{8}{3}{1}\rack{8}{3}{2}
				\rack{8}{3}{3}\rack{0}{2}{0}\rack{1}{2}{0}\rack{2}{2}{0}\rack{3}{2}{0}
				\rack{5}{2}{0}\rack{6}{2}{1}\rack{7}{2}{2}\rack{8}{2}{0}\rack{8}{2}{1}
				\rack{8}{2}{2}\obs{0}{1}{1}\obs{0}{1}{2}\obs{0}{1}{3}\obs{0}{1}{4}
				\obs{0}{1}{5}\obs{1}{1}{1}\portal{1}{1}{2}\portal{1}{1}{3}
				\portal{1}{1}{4}\obs{1}{1}{5}\obs{2}{1}{1}\portal{2}{1}{2}
				\portal{2}{1}{3}\portal{2}{1}{4}\obs{2}{1}{5}\rack{3}{1}{0}
				\obs{3}{1}{1}\obs{3}{1}{2}\obs{3}{1}{3}\obs{3}{1}{4}\obs{3}{1}{5}
				\rack{5}{1}{0}\rack{6}{1}{1}\rack{7}{1}{1}\rack{8}{1}{0}\rack{8}{1}{1}
				\rack{0}{0}{0}\rack{1}{0}{0}\rack{2}{0}{0}\rack{3}{0}{0}\rack{5}{0}{0}
				\rack{6}{0}{0}\rack{7}{0}{0}\rack{8}{0}{0}
		\end{scope}
		% the Overworld
		\chunkblock{2.4}
		\begin{scope}[voxels=2.4]
			\grass{5}{8}{0}\grass{6}{8}{1}\grass{7}{8}{0}\grass{0}{7}{0}
				\grass{1}{7}{1}\grass{2}{7}{1}\grass{3}{7}{0}\grass{5}{7}{0}
				\grass{6}{7}{0}\grass{7}{7}{0}\grass{0}{6}{0}\grass{1}{6}{1}
				\grass{2}{6}{1}\grass{3}{6}{0}\grass{5}{6}{0}\grass{6}{6}{0}
				\grass{7}{6}{0}\grass{0}{5}{0}\grass{1}{5}{0}\grass{2}{5}{0}
				\grass{3}{5}{0}\leaves{5}{4}{4}\leaves{6}{4}{5}\leaves{7}{4}{4}
				\grass{5}{3}{1}\leaves{5}{3}{5}\dirt{6}{3}{0}\grass{6}{3}{1}
				\wood{6}{3}{2}\wood{6}{3}{3}\leaves{6}{3}{5}\leaves{7}{3}{4}
				\leaves{7}{3}{5}\grass{0}{2}{0}\grass{1}{2}{0}\grass{2}{2}{0}
				\grass{3}{2}{0}\grass{5}{2}{1}\leaves{5}{2}{4}\dirt{6}{2}{0}
				\grass{6}{2}{1}\leaves{6}{2}{4}\leaves{6}{2}{5}\leaves{7}{2}{4}
				\obs{0}{1}{1}\obs{0}{1}{2}\obs{0}{1}{3}\obs{0}{1}{4}\obs{0}{1}{5}
				\obs{1}{1}{1}\portal{1}{1}{2}\portal{1}{1}{3}\portal{1}{1}{4}
				\obs{1}{1}{5}\obs{2}{1}{1}\portal{2}{1}{2}\portal{2}{1}{3}
				\portal{2}{1}{4}\obs{2}{1}{5}\grass{3}{1}{0}\obs{3}{1}{1}\obs{3}{1}{2}
				\obs{3}{1}{3}\obs{3}{1}{4}\obs{3}{1}{5}\grass{5}{1}{1}\grass{6}{1}{1}
				\grass{7}{1}{1}\dirt{8}{1}{0}\grass{8}{1}{1}\grass{0}{0}{0}
				\grass{1}{0}{0}\grass{2}{0}{0}\grass{3}{0}{0}\dirt{5}{0}{0}
				\grass{5}{0}{1}\dirt{6}{0}{0}\grass{6}{0}{1}\dirt{7}{0}{0}
				\grass{7}{0}{1}\dirt{8}{0}{0}\grass{8}{0}{1}\chest{8}{0}{2}
		\end{scope}
		% the two portals are glued along a common face
		\draw[wormhole] (0.5,0.2,3.125+0.025) to[out=180, in=90] (-0.75,0.2,1.875);
		\draw[wormhole] (-0.75,0.2,1.875) to[out=-90, in=180] (0.5,0.2,0.725+0.025);
		\draw[cellglyph] (-0.75,0.1,1.775) -- ++(0,0.2,0) -- ++(0,0,0.2) -- ++(0,-0.2,0) -- cycle;
		\node[anchor=west, font=\small\scshape] at (2.55,1,3.0) {Overworld};
		\node[anchor=west, font=\small\scshape] at (2.55,1,0.6) {Nether};
	\end{tikzpicture}
	\caption{A block sheaf over $Y$. The chunk at the front left of the
		Overworld and the chunk at the front left of the Nether each contain a
		Nether portal, and the wormhole face of $Y$ glues the two portals
		together.}
	\label{fig:nether-portal}
\end{figure}
